\section{Experiments}

\subsection{Experimental Questions}
The empirical section is organized around four questions:
\begin{enumerate}
    \item Do existing evaluators lose discriminative power as subject count increases?
    \item Does AnchorScore better align with pairwise human preference than standard image-text or image-image metrics?
    \item Does multi-task diagnostic supervision improve over ranking-only training?
    \item How stable are the conclusions across backbone scales and update modes?
\end{enumerate}

\subsection{Ready-To-Run Experimental Setup}
The current experiment plan is already fixed even though the final tables are still being collected. We train six AnchorScore variants using Qwen3.5 backbones at \texttt{0.8B}, \texttt{2B}, and \texttt{4B} scales, each under \texttt{layer\_only} and \texttt{lora\_layer} modes. All six runs use the same paper-facing budget of 600 optimizer steps. This shared-budget design is intended to make the main comparison reviewer-friendly under a realistic compute constraint. Training is performed on the 60k-pair dual-teacher silver pipeline, while paper-facing comparison is centered on the broader 4k-pair human evaluation set.

\subsection{Baselines}
The current baseline panel for metric comparison includes the evaluators that are already integrated into our analysis workflow: CLIP-T, CLIP-I, DINOv2, ArcFace, ImageReward, SCR, and REFVNL. These methods were selected because together they cover text-image similarity, image-image similarity, identity-oriented features, reward-model style preference signals, and specialized reference-aware scoring. In the final paper version, these baselines will be evaluated on the same multi-generator human set used to assess AnchorScore.

\subsection{What Can Be Written Now}
Even before the final result tables are complete, the following parts of the section are ready to write and should remain in the paper draft:
\begin{itemize}
    \item the evaluation questions and success criteria;
    \item the fixed training-budget protocol for the six AnchorScore variants;
    \item the distinction between the 60k-pair silver training set and the 4k-pair human evaluation set;
    \item the baseline list and why each baseline is included;
    \item the planned breakdown by subject count $N \in \{2,4,6,8\}$;
    \item the list of ablations and qualitative analyses to report.
\end{itemize}

\subsection{Evaluation Data Regime}
Our evaluation design intentionally separates scale from trustworthiness. The silver set supplies the large number of pairwise comparisons needed to train a reference-conditioned evaluator, but the central empirical claim of the paper is assessed on human judgments. The 4k-pair human evaluation set also spans more generator families than the training set, which makes it possible to test cross-generator transfer rather than only in-domain agreement on Nano and Mosaic.

\subsection{What Should Stay Empty For Now}
The following subsections should remain explicit placeholders until the full runs and evaluation outputs are finalized:
\begin{itemize}
    \item \textbf{Main result table}: final metric-to-human alignment numbers;
    \item \textbf{Subject-count scaling}: plots or tables showing score separability as $N$ increases;
    \item \textbf{Backbone comparison}: the final comparison across \texttt{0.8B}, \texttt{2B}, and \texttt{4B};
    \item \textbf{Mode comparison}: the final comparison between \texttt{layer\_only} and \texttt{lora\_layer};
    \item \textbf{Qualitative case studies}: curated failure and success examples with figures.
\end{itemize}

\subsection{Planned Main Results Table}
Table~\ref{tab:main_results_placeholder} is reserved for the final paper-facing summary once all runs finish.

\begin{table}[t]
    \centering
    \small
    \caption{Placeholder for the main result table. Fill after all evaluation outputs are finalized.}
    \label{tab:main_results_placeholder}
    \resizebox{\linewidth}{!}{%
    \begin{tabular}{lcccc}
        \toprule
        Method & Pair Acc. & Spearman & Kendall-$\tau$ & Notes \\
        \midrule
        CLIP-T & -- & -- & -- & text-image baseline \\
        CLIP-I & -- & -- & -- & image-image baseline \\
        DINOv2 & -- & -- & -- & self-supervised visual baseline \\
        ArcFace & -- & -- & -- & identity baseline \\
        ImageReward & -- & -- & -- & reward-model baseline \\
        SCR & -- & -- & -- & subject completeness baseline \\
        REFVNL & -- & -- & -- & reference-aware baseline \\
        \midrule
        AnchorScore (0.8B, layer\_only) & -- & -- & -- & pending \\
        AnchorScore (0.8B, lora\_layer) & -- & -- & -- & pending \\
        AnchorScore (2B, layer\_only) & -- & -- & -- & pending \\
        AnchorScore (2B, lora\_layer) & -- & -- & -- & pending \\
        AnchorScore (4B, layer\_only) & -- & -- & -- & pending \\
        AnchorScore (4B, lora\_layer) & -- & -- & -- & pending \\
        \bottomrule
    \end{tabular}
    }
\end{table}

\subsection{Planned Ablations and Analyses}
Once the runs and evaluation files are complete, we will populate this section with:
\begin{itemize}
    \item \textbf{Meta-evaluation against human judgment}: Kendall-$\tau$, Spearman, and pairwise accuracy on the held-out human set.
    \item \textbf{High-subject-count separability}: score-gap and overlap analysis between stronger and weaker generations for $N \in \{2,4,6,8\}$.
    \item \textbf{Cross-generator transfer}: comparison between generators seen in silver training data and generators appearing only in human evaluation.
    \item \textbf{Ablations}: ranking-only versus multi-task training, with and without agreement-based filtering.
    \item \textbf{Qualitative case studies}: examples where global metrics collapse but AnchorScore preserves the correct ordering.
\end{itemize}

The central experimental claim to validate remains the same: existing metrics lose score separability as subject count increases, whereas a reference-conditioned evaluator trained with pairwise ranking and diagnostic supervision can remain more discriminative and better aligned with human preference.
